\begin{courseunit}{III.1}
\chapter{The Independent Missal Reader}

\begin{learningobjectives}
  \item define a testable meaning of ``without a facing translation'';
  \item read a Missal passage in grammatical passes instead of guessing from memory;
  \item reconstruct prose order without erasing emphasis;
  \item use a dictionary only after morphology and clause structure have narrowed the question.
\end{learningobjectives}

Volume III assumes that the paradigms and syntax of the first two volumes are
available from memory. Its task is different: to make those resources operate
together on the page of the \Missal. A reader who can recite a familiar prayer
may still be unable to construe an unfamiliar collect. A reader who can produce
an elegant English version may still be unable to say which Latin word governs
which. The standard here is therefore visible analysis followed by direct
understanding.

For this volume, a \Term{complete Mass} means the fixed prayer text of the
\emph{Ordo Missæ} and Roman Canon together with every prayer, lesson, Gospel,
and chant text assigned to a particular formulary. It does not mean memorizing
rubrics, reading musical notation, or translating the priest's optional prayers
before and after Mass. The \emph{Kyrie} is Greek; Hebrew or Greek liturgical
words such as \emph{Amen}, \emph{Alleluia}, and \emph{Hosanna} remain part of
the corpus but are not evidence about Latin grammar.

Quoted Missal text follows the 1962 typical edition's house typography:
consonantal \Latin{I} in \Latin{Iesus, eius, iustitia}, and the printed ligature
\Latin{æ}. Other learner editions may expand \Latin{æ} to \Latin{ae} or print
consonantal \Latin{j}; those are orthographic variants, not new lemmas. Preserve
the source's spelling when copying, but search a lexicon under its normalized
headword when necessary.

\begin{grammarbox}{Working definitions for independent reading}
A \Term{genre} is a recurring class of text identified by its liturgical job
and its usual formal features: a Collect and a Gospel may concern the same
feast, but they remain different genres because one petitions and the other
proclaims a reading. A \Term{formulary} is the set of texts assigned together
for one Mass; it can therefore contain several genres.

A \Term{lemma} is the dictionary headword under which an inflected form is
listed. \Term{Morphology} identifies what form an ending or stem makes
(for example, first-person plural present subjunctive); \Term{syntax}
identifies the form's relation to other words (for example, the finite verb of
a purpose clause). A \Term{finite verb} is marked for person and number and
can anchor a clause; an infinitive, participle, gerund, or gerundive is
\Term{nonfinite} and depends on some larger construction.

A \Term{head} is the central word that a phrase expands. A \Term{governor} is
a word or construction that licenses another element, often demanding a case
or nonfinite form. A \Term{complement} completes the sense or construction of
its governor; a \Term{modifier} adds description without filling such a
required slot. An \Term{antecedent} is the expression to which a pronoun or
relative refers. \Term{Ellipsis} is the omission of wording recoverable from
grammar or strict context; it is not permission to insert whatever makes an
English version smooth.

To \Term{construe} a passage is to account for these forms and dependencies so
that its sense follows from the Latin. A \Term{prose-order reconstruction} is
a temporary diagnostic rearrangement that places dependencies together; it is
neither a correction of the source nor a translation.
\end{grammarbox}

\section{What ``no facing translation'' does and does not mean}

It does not forbid English during instruction. A precise translation can expose
a false parse. It does forbid using a parallel English column as a substitute
for reading during a gate. At a no-facing-translation gate the Latin page is
alone. The learner demonstrates understanding by marking clauses, identifying
forms, answering short Latin questions, giving a restrained written sense, and
explaining the relation of each subordinate clause to its governor.

Familiarity is not the same as sight reading. Accordingly, every final gate has
two parts: a familiar text read fluently and an unseen text of the same genre.
\Term{Unseen} means not used in the teaching volume or assigned worksheet,
although its grammar and high-frequency vocabulary may be familiar.

\begin{grammarbox}{Three passes through an unfamiliar passage}
\textbf{Pass 1---frame.} Circle coordinators and subordinators; underline every
finite verb; bracket vocatives and parenthetical \Latin{quæsumus}; draw a line
after each independent clause. Do not use a dictionary.

\textbf{Pass 2---dependencies.} Mark subjects, direct objects, governed cases,
infinitives, participles, and antecedents. Supply an omitted form only in square
brackets. Write the supplied Latin, not an English guess.

\textbf{Pass 3---sense.} Read again in the printed order. Consult a dictionary
only for a word whose lemma and syntactic slot have already been proposed.
Revise the clause map when the proposed meaning does not fit the construction.
\end{grammarbox}

\section{Reconstruction without flattening}

Missal word order is meaningful even when it is not the order most convenient
for first analysis. Make a disposable prose-order line beneath the text:

\begin{enumerate}
  \item write the finite verb and its explicit or implicit subject;
  \item attach complements demanded by that verb;
  \item place modifiers beside their heads;
  \item expand subordinate clauses from the outside inward;
  \item return to the printed order and state what the displacement emphasizes.
\end{enumerate}

Consider the controlled drill \Latin{Tibi, Domine, sacrificium laudis
offerimus}, built from recurrent Missal offering diction but not presented as a
quotation. Analytical order is
\Latin{nos offerimus sacrificium laudis tibi, Domine}. The initial
\Latin{tibi} is not disorder: it foregrounds the recipient. Reconstruction is
a scaffold, not a rewritten liturgical text.

\section{Dictionary discipline}

Never look up an inflected form as though it were a lemma. Before opening a
dictionary, write four things: proposed lemma, morphology, syntactic function,
and a provisional semantic class. For \Latin{perfruamur}, for example, write
\Latin{perfruor, perfrui, perfructus sum; first plural present subjunctive,
deponent; verb of the ut-clause; enjoyment/possession; governs ablative}.

During unit work, allow at most five new content-word lookups per one hundred
words on the first assisted reading. Function words, endings, and words already
entered in the vocabulary ledger do not earn repeat lookups. Record a family,
not an orphan: \Latin{redimo, redemptio, redemptor}; \Latin{offero, oblatio,
munus}; \Latin{salvus, salvo, salus, salutaris}.

\MasteryGate{on an unseen 120--180-word Missal passage, you can mark every
finite clause with at least 90\% accuracy, identify the lemma and morphology of
at least 90\% of sampled forms, justify every supplied ellipsis, and state the
literal sense with no facing text and no more than five first-pass lookups.}


\end{courseunit}

\begin{courseunit}{III.2}
\chapter{The Ordinary and the Fixed Order}

\begin{learningobjectives}
  \item distinguish the sung Ordinary from the wider fixed \emph{Ordo Missæ};
  \item resolve apposition, relative clauses, and compact liturgical responses;
  \item read familiar fixed texts by grammar rather than by recited memory.
\end{learningobjectives}

For this course, the recurring sung Ordinary corpus comprises the
\emph{Kyrie}, \emph{Gloria}, \emph{Credo}, \emph{Sanctus--Benedictus}, and
\emph{Agnus Dei}. The fixed Order is larger: it
includes dialogues, prayers at the altar, offertory prayers, preparation for
Communion, and dismissal. A claim to read the ``whole Ordinary'' must say which
of these two corpora it means. This volume requires both, together with the
Roman Canon studied later.

Here \Term{Ordinary} means the recurring sung set just named, whereas the
\Term{Ordo Missæ} means the fixed ritual order and its texts. A \Term{corpus}
is simply a defined body of texts gathered for study. In either corpus,
\Term{apposition} places two expressions in the same syntactic role so that the
second identifies or redescribes the first; apposition is not a second,
unstated clause.

\section{Responses and legitimate ellipsis}

Recall that an \Term{ellipsis} omits words that grammar and the immediate
formula make recoverable. A completion is \Term{legitimate} only when it
preserves the construction and force of the printed words. An
\Term{optative} expression presents a wish; a statement presents something as
fact. That contrast determines which verb may be supplied in a response.

\Latin{Dominus vobiscum. Et cum spiritu tuo} (\Missal, \emph{Ordo Missæ})
contains no expressed verb. The exchange is conventionally construed
optatively; a pedagogical paraphrase of the response is \Latin{[Dominus sit]
cum spiritu tuo}. Because \Latin{sit} is not expressed or uniquely forced by
the morphology, mark it interpretive rather than grammatically required. The
preposition governs the ablative; \Latin{tuo} agrees with \Latin{spiritu}.
Supplying \Latin{est} would not preserve the conventional optative force.

In \Latin{Gratias agamus Domino Deo nostro. Dignum et iustum est} (ibid.),
\Latin{agamus} is a hortatory subjunctive, \Latin{Domino Deo} is dative after
\Latin{gratias agere}, and \Latin{nostro} modifies the appositional whole. The
reply has predicate neuter adjectives with an implicit infinitive or action:
``doing so is worthy and just.'' Do not invent a nominative masculine noun.

\section{Agreement across distance}

The Gloria opens, \Latin{Gloria in excelsis Deo, et in terra pax hominibus
bonæ voluntatis} (\emph{Gloria in excelsis}; compare Lk.~2:14). The balanced
prepositional phrases locate two nominative ideas, \Latin{gloria} and
\Latin{pax}; \Latin{Deo} and \Latin{hominibus} are datives. In the Missal text
\Latin{bonæ voluntatis} is genitive singular and qualifies the people: human
beings characterized by good will. It is not nominative with \Latin{pax}.

The Creed says, \Latin{Qui propter nos homines et propter nostram salutem
descendit de cælis} (\emph{Symbolum Nicæno-Constantinopolitanum}).
\Latin{homines} is accusative in apposition to \Latin{nos}; the repeated
\Latin{propter} keeps persons and purpose parallel. The relative \Latin{qui}
continues the confessed subject from the preceding clause.

\section{Purpose clauses and prayer diction}

At the ascent to the altar the priest says:

\begin{quote}
\Latin{Aufer a nobis, quæsumus, Domine, iniquitates nostras: ut ad Sancta
sanctorum puris mereamur mentibus introire.}

\emph{Source:} \Missal, \emph{Ordo Missæ}, prayer \emph{Aufer a nobis}.
\end{quote}

The imperative \Latin{aufer} governs both the object \Latin{iniquitates
nostras} and separation phrase \Latin{a nobis}. The goal is expressed by
\Latin{ut} with \Latin{mereamur}; \Latin{introire} complements that deponent.
\Latin{ad Sancta sanctorum} is biblical genitive intensification: accusative
neuter plural \Latin{Sancta} plus genitive neuter plural \Latin{sanctorum},
conventionally ``to the Holy of Holies.'' \Latin{puris mentibus} is ablative of
manner or attendant state.

The offertory prayer begins:

\begin{quote}
\Latin{Suscipe, sancta Trinitas, hanc oblationem, quam tibi offerimus ob
memoriam passionis, resurrectionis et ascensionis Iesu Christi Domini nostri.}

\emph{Source:} \Missal, \emph{Ordo Missæ}, \emph{Suscipe, sancta Trinitas}.
\end{quote}

\Latin{Trinitas} is vocative in function although formally nominative;
\Latin{quam} is the direct object of \Latin{offerimus} and refers to
\Latin{oblationem}. The dative \Latin{tibi} precedes its verb. The preposition
\Latin{ob} governs \Latin{memoriam}; the three dependent genitives name what is
remembered, and \Latin{Iesu Christi Domini nostri} depends on all three.

\section{Vocabulary families of the fixed Order}

\begin{description}[style=nextline]
  \item[Offering:] \Latin{offero, offerre, obtuli, oblatum}; \Latin{oblatio};
    \Latin{munus}; \Latin{hostia}; \Latin{sacrificium}.
  \item[Receiving:] \Latin{accipio}; \Latin{suscipio}; \Latin{receptus};
    \Latin{perceptio}. Prefixes change viewpoint, not merely English wording.
  \item[Holiness:] \Latin{sanctus}; \Latin{sanctifico}; \Latin{sacratus};
    \Latin{sacrosanctus}.
  \item[Mercy and worthiness:] \Latin{misericors, misericordia, misereor};
    \Latin{dignus, dignor, indignus}.
  \item[Faith and confession:] \Latin{credo, fides, fidelis};
    \Latin{confiteor, confessio}.
\end{description}

\MasteryGate{you can read every congregational fixed text and the assigned
priestly fixed prayers silently, on the page, without a facing translation; on
twenty randomly selected forms you score at least 18/20 for morphology and
function; and you give a direct written sense of ten short responses with at
least 95\% content accuracy.}


\end{courseunit}

\begin{courseunit}{III.3}
\chapter{Collects: The Compressed Petition}

\begin{learningobjectives}
  \item identify the address, amplification, petition, desired result, and conclusion of a collect;
  \item untangle nested relatives, ablative absolutes, and paired gerunds;
  \item preserve the rhetorical focus after making a prose-order reconstruction.
\end{learningobjectives}

A Roman collect is a compact period, not a string of devotional slogans. Its
usual architecture is (1) divine address, (2) a relative or participial reason
for confidence, (3) one governing petition, (4) an \Latin{ut} or \Latin{ne}
clause stating the desired effect, and (5) a mediation conclusion. Not every
collect contains every member, and punctuation is a reading aid rather than a
substitute for syntax.

A \Term{Collect} is the presiding oration that gathers the assembly's petition
into a compact prayer. A \Term{period} is a sentence designed as an articulated
whole whose subordinate members lead toward or depend on a governing member;
it does not merely mean a long sentence. The \Term{address} names the person
addressed; an \Term{amplification} expands that address, commonly by recalling
a divine attribute or act; the \Term{petition} states what is requested; the
\Term{desired result} states the end for which it is requested; and the
\Term{conclusion} supplies the conventional mediation. These are functional
members, not guaranteed punctuation units.

\section{Imperatives, infinitives, and ablative absolutes}

An \Term{imperative} directly enjoins an addressee. An \Term{infinitive} is a
nonfinite verbal form that can fill noun-like roles; its subject, when expressed
in its own clause, is normally accusative. An \Term{ablative absolute} is a
grammatically detached ablative nominal with an agreeing participle or other
predicate; it supplies attendant circumstance without serving as the subject,
object, or other required complement of the main verb.

\begin{quote}
\Latin{Excita, quæsumus, Domine, potentiam tuam, et veni: ut ab imminentibus
peccatorum nostrorum periculis, te mereamur protegente eripi, te liberante
salvari.}

\emph{Source:} \Missal, Dominica I Adventus, Collecta.
\end{quote}

First isolate \Latin{excita} and \Latin{veni}, two coordinated imperatives.
\Latin{quæsumus} is parenthetical and \Latin{Domine} vocative. Inside the
\Latin{ut}-clause, \Latin{mereamur} governs the passive infinitives
\Latin{eripi} and \Latin{salvari}. The separation phrase
\Latin{ab imminentibus ... periculis} belongs especially with \Latin{eripi}.
Both \Latin{te protegente} and \Latin{te liberante} are ablative absolutes; the
pronoun is their subject, not the object of the infinitives.

Analytical order:
\begin{quote}
\Latin{Domine, excita potentiam tuam et veni, ut mereamur eripi ab periculis
peccatorum nostrorum te protegente et salvari te liberante.}
\end{quote}
Return then to the printed order: danger is placed before protection, and the
balanced ablative absolutes make divine action surround the two infinitives.

\section{Correlative relatives}

A \Term{correlative} pair uses matched words to link corresponding roles. In
the genitive pair \Latin{cuius ... eius}, the first member introduces the
relative relation (``whose'') and the second resumes it (``of that same
one''); agreement and case still have to be parsed separately.

\begin{quote}
\Latin{Deus, qui hanc sacratissimam noctem veri luminis fecisti illustratione
clarescere: da, quæsumus; ut, cuius lucis mysteria in terra cognovimus, eius
quoque gaudiis in cælo perfruamur.}

\emph{Source:} \Missal, In Nativitate Domini, ad primam Missam in nocte,
Collecta.
\end{quote}

The relative \Latin{qui} addresses the same \Latin{Deus}; its second-person
verb \Latin{fecisti} is normal because the antecedent is vocative in function.
\Latin{facere} with infinitive appears as \Latin{fecisti ... clarescere};
\Latin{illustratione} is ablative of means, and \Latin{veri luminis} defines
the illumination. The petition is the bare imperative \Latin{da}; its content
is the following \Latin{ut}-clause. \Latin{cuius ... eius} is a correlative:
the light whose mysteries were known is the light whose joys are enjoyed.
\Latin{perfruamur} is deponent and governs ablative \Latin{gaudiis}.

\section{Gerunds as divine accompaniment}

A \Term{gerund} is a neuter singular verbal noun, active in meaning, used only
in oblique cases. It differs from a \Term{gerundive}, which is a verbal
adjective agreeing with a noun and often expressing necessity. Thus the
ablatives \Latin{præveniendo} and \Latin{adiuvando} name ways of acting; they
do not agree with \Latin{vota}.

\begin{quote}
\Latin{Deus, qui hodierna die per Unigenitum tuum æternitatis nobis aditum,
devicta morte, reserasti: vota nostra, quæ præveniendo aspiras, etiam
adiuvando prosequere.}

\emph{Source:} \Missal, Dominica Resurrectionis, Collecta.
\end{quote}

The main petition is \Latin{prosequere}, a deponent imperative, with
\Latin{vota nostra} as object. The relative \Latin{quæ} also has
\Latin{vota} as antecedent and object of \Latin{aspiras}. The ablative gerunds
\Latin{præveniendo} and \Latin{adiuvando} express the manner or accompanying
divine action: grace inspires desires beforehand and carries them through by
helping. In the relative amplification, \Latin{aditum} is the object of
\Latin{reserasti}, \Latin{nobis} dative of advantage, and
\Latin{devicta morte} an ablative absolute.

\section{A collect lexicon}

Learn families in their recurrent constructions:
\begin{itemize}
  \item \Latin{præsta, concede, tribue, largire, da} introduce petitions;
  \item \Latin{mereor + infinitive} means ``be found or deemed worthy to'';
    \Latin{dignor + infinitive}, ``deign or be pleased to''; and
    \Latin{valeam + infinitive}, ``be able to.'' Determine the wider theology
    from the prayer and its sources rather than from these constructions alone;
  \item \Latin{propitius} is often predicative/adverbial with an imperative;
  \item \Latin{familia, populus, Ecclesia, vota, desideria} commonly become
    subjects or objects far from their modifiers;
  \item \Latin{prævenio, prosequor, perficio, perago} describe grace beginning,
    accompanying, and completing an action.
\end{itemize}

\MasteryGate{on four unseen collects from different seasons, you can label
every architectural member, reconstruct each main and subordinate clause,
parse all nonfinite verbs, and give a literal sense with at least 90\% clause
accuracy and no facing translation; no collect may contain an unaccounted
finite verb.}


\end{courseunit}

\begin{courseunit}{III.4}
\chapter{Secrets and Postcommunions: Gift and Effect}

\begin{learningobjectives}
  \item recognize the ritual vocabulary and demonstratives of the Secret;
  \item follow sacramental cause and desired effect in the Postcommunion;
  \item distinguish the grammatical referents of \Latin{quæ}, \Latin{hæc}, and \Latin{quam}.
\end{learningobjectives}

The Secret often moves from gifts presently offered to the unity, cleansing,
or salvation they signify. The Postcommunion moves from reception to continuing
effect. Both genres favor passive infinitives, optative subjunctives, abstract
nouns, and pronouns whose antecedents must be found by gender and number rather
than proximity.

In the 1962 books, the \Term{Secret} is the proper oration said silently over
the offerings before the Preface; the \Term{Postcommunion} is the proper
oration after Communion that asks for the received sacrament's continuing
effect. These names identify liturgical position and function, not a particular
sentence pattern. A \Term{referent} is the person, thing, or idea to which an
expression points; its \Term{antecedent} is the earlier wording that supplies
that referent. Antecedent matching begins with gender, number, and syntactic
possibility, then tests whether the resulting sense fits the prayer.

An \Term{abstract noun} names a quality, act, or state---such as unity,
reception, or cleansing---rather than a concrete bearer. A \Term{sacramental
effect} is the reality the prayer asks the received or offered sacrament to
bring about; it must be distinguished from the ritual gift, sign, or act that
the sentence makes its grammatical cause or occasion.

\section{The offered sign}

\begin{quote}
\Latin{Ecclesiæ tuæ, quæsumus, Domine, unitatis et pacis propitius dona
concede: quæ sub oblatis muneribus mystice designantur.}

\emph{Source:} \Missal, In festo Sanctissimi Corporis Christi, Secreta.
\end{quote}

The main object is \Latin{dona}; \Latin{unitatis et pacis} specifies the gifts,
and \Latin{Ecclesiæ tuæ} is dative of recipient. Read \Latin{propitius} with
the subject implicit in \Latin{concede}. Neuter plural \Latin{quæ} refers to
\Latin{dona}, not to feminine \Latin{Ecclesiæ} or the singular genitives.
\Latin{sub oblatis muneribus} presents the gifts under the outward offered
elements; \Latin{designantur} is passive.

\section{Reception as a sign of future fruition}

\begin{quote}
\Latin{Fac nos, quæsumus, Domine, divinitatis tuæ sempiterna fruitione
repleri: quam pretiosi Corporis et Sanguinis tui temporalis perceptio
præfigurat.}

\emph{Source:} \Missal, In festo Sanctissimi Corporis Christi, Postcommunio.
\end{quote}

\Latin{fac} takes accusative \Latin{nos} and passive infinitive
\Latin{repleri}. The ablative \Latin{fruitione} expresses that with which the
recipient is filled; \Latin{divinitatis tuæ} is its objective genitive.
Feminine singular \Latin{quam} refers to \Latin{fruitionem} understood from
\Latin{fruitione}. The subject of \Latin{præfigurat} is \Latin{temporalis
perceptio}; the genitives \Latin{pretiosi Corporis et Sanguinis tui} specify
what is received.

\section{A future feast approached now}

\begin{quote}
\Latin{Suscipiamus, Domine, misericordiam tuam in medio templi tui: ut
reparationis nostræ ventura solemnia congruis honoribus præcedamus.}

\emph{Source:} \Missal, Dominica I Adventus, Postcommunio.
\end{quote}

\Latin{suscipiamus} is hortatory or petitionary subjunctive. In the result or
purpose clause, \Latin{ventura solemnia} is the accusative object of
\Latin{præcedamus}: the worshippers go before the coming solemnities with
fitting honors. \Latin{reparationis nostræ} defines those solemnities and
\Latin{congruis honoribus} is ablative of means or manner. Do not make
\Latin{solemnia} the subject merely because its form could be nominative.

\section{Semantic families}

\begin{description}[style=nextline]
  \item[Gift:] \Latin{munus} emphasizes a gift or duty; \Latin{donum} a gift;
    \Latin{oblatio} the offering; \Latin{hostia} the sacrificial victim.
  \item[Reception:] \Latin{percipio/perceptio} is full reception or taking;
    \Latin{sumo/sumptio} taking up; \Latin{suscipio} receiving or sustaining.
  \item[Effect:] \Latin{purgo, mundo, emundo}; \Latin{reficio};
    \Latin{reparo}; \Latin{proficiat}; \Latin{medela, remedium, fructus}.
  \item[Sign:] \Latin{significo}; \Latin{designo}; \Latin{præfiguro};
    \Latin{mysterium}; \Latin{sacramentum}.
\end{description}

\MasteryGate{given two unseen Secrets and two Postcommunions, you identify the
offered or received reality, every requested effect, and every pronominal
antecedent with at least 90\% accuracy; you may use a Latin dictionary after a
complete first-pass clause map but no facing translation.}


\end{courseunit}

\begin{courseunit}{III.5}
\chapter{Lessons and Epistles: Argument on the Page}

\begin{learningobjectives}
  \item map argument markers and pronoun reference in Vulgate prose;
  \item construe indirect statement, participial chains, and biblical idiom;
  \item summarize a lesson's logical movement before translating individual words.
\end{learningobjectives}

The liturgical heading \Latin{Lectio Epistolæ beati Pauli Apostoli} is not part
of Paul's sentence. Likewise \Latin{Fratres} is a liturgical vocative added or
retained at the reading's opening. Mark such frames before tracing an argument.
Then circle \Latin{enim, autem, ergo, itaque, igitur, quoniam, sicut}, and
\Latin{sed}. These small words reveal the skeleton that freer word order hides.

A \Term{Lesson} is a proclaimed scriptural reading; an \Term{Epistle} is a
Lesson presented under an epistolary heading or drawn from an apostolic
letter. The label does not change the quoted passage's grammar. An
\Term{argument} is an ordered movement from grounds through relations such as
contrast or inference to a claim or exhortation. An \Term{argument connector}
is a word that signals one of those relations: a reason supports a claim, a
contrast sets opposed members against each other, and an inference presents a
conclusion as following from what precedes.

\section{Exhortation built on time}

\begin{quote}
\Latin{Fratres: Scientes quia hora est iam nos de somno surgere. Nunc enim
propior est nostra salus, quam cum credidimus. Nox præcessit, dies autem
appropinquavit. Abiciamus ergo opera tenebrarum, et induamur arma lucis.}

\emph{Source:} \Missal, Dominica I Adventus, Epistola (Rom.~13:11--12).
\end{quote}

\Latin{scientes} agrees with the implied plural addressees and introduces the
reason for the exhortation. After \Latin{hora est}, the accusative-infinitive
phrase \Latin{nos ... surgere} states what it is time for. Comparative
\Latin{propior} agrees with \Latin{salus}; \Latin{quam cum credidimus} compares
the present with the time when belief began. The paired perfects establish the
new situation; \Latin{abiciamus} and passive \Latin{induamur} are hortatory
subjunctives. \Latin{opera tenebrarum} and \Latin{arma lucis} make the ethical
contrast concrete.

\section{Tradition received and handed on}

A \Term{participial phrase} consists of a participle with its complements and
modifiers. Because a participle is verbal but agrees adjectivally with a noun
or pronoun, its agreement identifies its participant while its tense normally
expresses time relative to the governing verb, not an independent absolute
date.

\begin{quote}
\Latin{Ego enim accepi a Domino quod et tradidi vobis: quoniam Dominus Iesus,
in qua nocte tradebatur, accepit panem, et gratias agens fregit.}

\emph{Source:} \Missal, In Cena Domini and In festo Sanctissimi Corporis
Christi, Epistola (1~Cor.~11:23--24).
\end{quote}

The neuter relative \Latin{quod} is simultaneously what Paul received and what
he handed on. The dative \Latin{vobis} belongs to \Latin{tradidi}. In
\Latin{in qua nocte}, \Latin{qua} is a relative adjective agreeing with
\Latin{nocte}; no separate antecedent is printed: ``on which night.'' Imperfect passive
\Latin{tradebatur} and perfect active \Latin{accepit} deliberately place
betrayal and gift together. \Latin{gratias agens} is a present participial
phrase whose action accompanies and precedes \Latin{fregit}.

\section{A difficult indirect statement}

An \Term{indirect statement} reports a thought, perception, or assertion as a
dependent accusative-and-infinitive construction: the reported subject is
accusative and its verb infinitive. A predicate noun or adjective within that
construction normally agrees with the accusative subject. Do not confuse this
with direct quotation, which preserves the speaker's finite wording.

\begin{quote}
\Latin{Hoc enim sentite in vobis, quod et in Christo Iesu: qui cum in forma Dei
esset, non rapinam arbitratus est esse se æqualem Deo: sed semetipsum
exinanivit, formam servi accipiens.}

\emph{Source:} \Missal, Dominica II Passionis seu in Palmis, Epistola
(Phil.~2:5--7).
\end{quote}

The imperative \Latin{sentite} asks for the disposition that is also in Christ;
\Latin{quod} supplies the comparison's content. Concessive or circumstantial
\Latin{cum ... esset} frames the main contrast. With \Latin{arbitratus est},
the indirect statement is \Latin{se esse æqualem Deo}: \Latin{se} is its
accusative subject, \Latin{æqualem} predicate accusative, and \Latin{Deo}
dative after \Latin{æqualis}. \Latin{rapinam} is the predicate complement of
what he did not judge equality to be. The adversative \Latin{sed} turns to
\Latin{semetipsum exinanivit}; \Latin{formam} is the object of participle
\Latin{accipiens}.

\section{Argument-map notation}

In the margin write \Term{R} for reason, \Term{C} for contrast, \Term{I} for
inference, and \Term{E} for exhortation. A useful summary is not yet a
translation: ``R: the hour has arrived; C: night/day; I/E: discard/put on.''
This prevents a long English sentence from concealing the Latin logic.

An \Term{argument map} records relations among propositions without replacing
their wording. A \Term{proposition} is the content asserted, denied, commanded,
or entertained by a clause; several clauses may cooperate in one step of an
argument, and one sentence may contain several propositions.

\MasteryGate{on an unseen 200--300-word Epistle, you mark at least 90\% of
finite clauses and argument connectors, correctly resolve all sampled pronouns
and indirect statements, and produce a five-step argument map plus a literal
sense without a facing translation and with no more than eight lookups.}


\end{courseunit}

\begin{courseunit}{III.6}
\chapter{Gospels: Narrative, Speech, and Biblical Idiom}

\begin{learningobjectives}
  \item separate narrative frame, reported speech, and narrator comment;
  \item track participants through pronouns, participles, and shifts of subject;
  \item recognize common Semitic patterns without declaring them meaningless idioms.
\end{learningobjectives}

The heading \Latin{In illo tempore} gives the liturgical reading a frame but
does not connect syntactically to the first Gospel verb. Draw quotation bars in
the margin whenever speech begins and ends. Within a quotation, re-establish the
speaker and addressee before resolving \Latin{ego, tu, hic, ille}, or reflexive
\Latin{se}.

The \Term{narrative frame} locates or introduces the narrated event; the
\Term{narrator} is the discourse voice that presents events but is not thereby
a participant in them. \Term{Direct speech} reproduces a speaker's words as a
quotation, whereas \Term{reported speech} makes those words depend
grammatically on a reporting verb. A \Term{biblical idiom} is a recurrent
expression characteristic of scriptural Latin, sometimes reflecting Hebrew or
Greek usage; calling it an idiom describes established use, not absence of
grammar.

\section{Repetition as theological precision}

\begin{quote}
\Latin{Omnia per ipsum facta sunt: et sine ipso factum est nihil, quod factum
est.}

\emph{Source:} \Missal, In Nativitate Domini, ad tertiam Missam in die,
Evangelium (Io.~1:3).
\end{quote}

Neuter plural \Latin{omnia} is the subject of plural \Latin{facta sunt};
\Latin{per ipsum} expresses agency or mediation. In the second member,
\Latin{nihil} is singular and therefore takes \Latin{factum est}. The final
relative \Latin{quod factum est} restricts the created totality. Preserve the
repetition of \Latin{factum}: replacing every occurrence with a different
English synonym obscures the argument.

\section{Impersonal frame and result}

\begin{quote}
\Latin{Factum est autem, cum essent ibi, impleti sunt dies ut pareret. Et
peperit filium suum primogenitum, et pannis eum involvit, et reclinavit eum in
præsepio.}

\emph{Source:} \Missal, In Nativitate Domini, ad primam Missam in nocte,
Evangelium (Lc.~2:6--7).
\end{quote}

\Latin{factum est} is an impersonal narrative hinge. The temporal
\Latin{cum}-clause takes subjunctive \Latin{essent}; the main event is
\Latin{impleti sunt dies}. The clause \Latin{ut pareret} is an epexegetic
result/content clause explaining what the completion of the days amounted to:
the time for her to give birth. Thereafter the repeated subject remains
Mary: \Latin{peperit, involvit, reclinavit}. The pronoun \Latin{eum} twice
returns to \Latin{filium}; \Latin{pannis} is ablative of wrapping or means and
\Latin{in præsepio} location.

\section{Participles and a moving subject}

\begin{quote}
\Latin{Cum natus esset Iesus in Bethlehem Iuda in diebus Herodis regis, ecce
Magi ab oriente venerunt Ierosolymam, dicentes: Ubi est qui natus est Rex
Iudæorum?}

\emph{Source:} \Missal, In Epiphania Domini, Evangelium (Mt.~2:1--2).
\end{quote}

The circumstantial \Latin{cum}-clause has \Latin{Iesus} as subject; the main subject
then shifts to \Latin{Magi}. \Latin{Ierosolymam} is accusative of motion toward
without a preposition. Present participle \Latin{dicentes} agrees with the Magi
and opens their direct speech. In the question, substantival \Latin{qui natus
est} is the subject and \Latin{Rex} its predicate nominative; the capital letters
do not determine syntax.

\section{The first day of the week}

\begin{quote}
\Latin{Una autem sabbati, Maria Magdalene venit mane, cum adhuc tenebræ
essent, ad monumentum: et vidit lapidem sublatum a monumento.}

\emph{Source:} \Missal, Sabbato in Albis, Evangelium (Io.~20:1).
\end{quote}

\Latin{una sabbati} is biblical idiom: ``one [day] of the week,'' hence the
first day. Supply \Latin{die} only as an explanatory ellipsis, not as though it
were printed. The temporal \Latin{cum}-clause takes \Latin{essent}. In
\Latin{lapidem sublatum}, the participle is predicative: she saw the stone in a
removed state; \Latin{a monumento} marks separation.

\MasteryGate{on an unseen Gospel of 250--400 words, you mark speech boundaries,
track every participant through sampled pronouns and participles, explain each
subjunctive, and answer eight factual and causal questions in simple Latin with
at least 85\% content accuracy, using no facing translation and at most ten
lookups after the first pass.}


\end{courseunit}

\begin{courseunit}{III.7}
\chapter{Graduals, Alleluias, Tracts, and Sequences}

\begin{learningobjectives}
  \item read psalmic parallelism without forcing every colon into prose syntax;
  \item restore only ellipses demanded by grammar;
  \item recognize poetic word order, alternate perfect endings, and predicative participles;
  \item paraphrase a chant unit in simple Latin before making an English version.
\end{learningobjectives}

Chant texts are short, but they are not necessarily easy. A Gradual or Tract
may excerpt distant psalm verses, omit a repeated subject, or coordinate clauses
by rhythm rather than conjunction. A sequence adds rhyme, compressed images,
and occasionally archaic forms. Treat the verse marker as a performance marker,
not always a syntactic boundary.

In this volume a \Term{chant unit} is the textual unit assigned to one chant
position, even when it draws on nonadjacent source verses. A \Term{Gradual} is
the chant after the Lesson, an \Term{Alleluia} is the alleluiatic chant before
the Gospel, and a \Term{Tract} takes the corresponding pre-Gospel place in
specified non-alleluiatic Masses. A \Term{Sequence} is an extended strophic
text placed before the Gospel in the Masses to which the 1962 Missal assigns
one, after the Alleluia or Tract as applicable. These liturgical definitions do
not imply that all four share one poetic syntax.

\Term{Parallelism} arranges corresponding members so that their likeness,
contrast, or progression is perceptible. A \Term{colon} is one such sense-and-
rhythm member of a verse; punctuation may mark it, but a colon need not be an
independent clause. \Term{Poetic word order} separates or reorders related
elements for pattern or emphasis while leaving their inflections and
dependencies intact.

\section{Psalmic parallelism}

\begin{quote}
\Latin{Universi, qui te exspectant, non confundentur, Domine. V. Vias tuas,
Domine, notas fac mihi: et semitas tuas edoce me.}

\emph{Source:} \Missal, Dominica I Adventus, Graduale (Ps.~24:3--4).
\end{quote}

Substantival \Latin{universi} is modified by \Latin{qui te exspectant} and is
subject of passive \Latin{confundentur}. The next verse balances
\Latin{vias tuas} with \Latin{semitas tuas}. In \Latin{notas fac mihi},
\Latin{notas} is predicate accusative with \Latin{vias}; \Latin{mihi} is dative.
The second imperative \Latin{edoce} takes \Latin{me} and may take the thing
taught as accusative. Do not supply a new divine subject between the parallel
members.

The Alleluia of the same Mass reads:
\begin{quote}
\Latin{Ostende nobis, Domine, misericordiam tuam: et salutare tuum da nobis.}

\emph{Source:} \Missal, Dominica I Adventus, Alleluia (Ps.~84:8).
\end{quote}
The two imperatives reverse object and recipient: \Latin{nobis ... misericordiam}
then \Latin{salutare ... nobis}. Neuter substantival \Latin{salutare} means the
saving help manifested by God; it is not an infinitive.

\section{The sustained Tract}

\begin{quote}
\Latin{Qui habitat in adiutorio Altissimi, in protectione Dei cæli
commorabitur. V. Dicet Domino: Susceptor meus es tu, et refugium meum.}

\emph{Source:} \Missal, Dominica I in Quadragesima, Tractus (Ps.~90:1--2).
\end{quote}

Substantival \Latin{qui} supplies the subject of future \Latin{commorabitur}.
The two \Latin{in} phrases use ablatives and form a parallel description of
dwelling. Future \Latin{dicet} takes dative \Latin{Domino}; the quoted sentence
has \Latin{tu} as subject and \Latin{susceptor meus} and \Latin{refugium meum}
as predicate nominatives. The long Tract continues by changing persons and
speakers; mark those changes rather than assuming a single uninterrupted voice.

\section{Sequence compression}

\begin{quote}
\Latin{Victimæ paschali laudes immolent Christiani. Agnus redemit oves:
Christus innocens Patri reconciliavit peccatores.}

\emph{Source:} \Missal, Dominica Resurrectionis, Sequentia
\emph{Victimæ paschali laudes}.
\end{quote}

\Latin{immolent} is jussive subjunctive; \Latin{Christiani} is its subject,
\Latin{laudes} its object, and \Latin{Victimæ paschali} dative recipient.
The next two clauses use stark nominative--verb--object patterns. Dative
\Latin{Patri} belongs with \Latin{reconciliavit}; \Latin{Christus innocens}
and the sacrificial \Latin{Agnus} have the same referent and form a theological
parallel, but they are the distinct subjects of \Latin{reconciliavit} and
\Latin{redemit}, respectively.

Later the sequence says:
\begin{quote}
\Latin{Mors et vita duello conflixere mirando: dux vitæ mortuus regnat vivus.}
\end{quote}
\Latin{conflixere} is third-person plural perfect with the alternative ending
\Latin{-ere} for \Latin{-erunt}; it is not a syncopated form. \Latin{duello
mirando} is ablative of the wondrous combat.
\Latin{mortuus} and \Latin{vivus} are predicative: the dead leader of life reigns
living. Prose order helps---\Latin{dux vitæ, mortuus, vivus regnat}---but the
crossed poetic order makes death and living meet in one line.

A \Term{third-person plural perfect in \Latin{-ere}} uses an inherited
alternative to the ending \Latin{-erunt}; the two endings have the same person,
number, tense, voice, and mood. A \Term{predicative} adjective or
participle states a condition of a subject or object within the verbal
assertion; an \Term{attributive} one merely identifies or describes its noun
inside the noun phrase. Thus \Latin{mortuus} and \Latin{vivus} tell how the
leader stands in relation to \Latin{regnat}, not just which leader is meant.

\section{A responsible ellipsis ledger}

For every supplied word write one of three labels:
\begin{description}[style=nextline]
  \item[G] grammatically required, as an implicit subject contained in a verb;
  \item[P] supplied from a strict parallel member;
  \item[I] interpretive, helpful for sense but not recoverable as exact wording.
\end{description}
Only G and carefully justified P supplies belong in a clause reconstruction.
An I supply may appear in commentary but never inside the Latin quotation.

\MasteryGate{on one unseen Gradual, one Tract section, and twelve unseen lines
of a sequence, you identify all finite forms, reconstruct prose order, label
every ellipsis G/P/I, and write a simple Latin paraphrase preserving the agents
and actions; the parse and dependency score is at least 90\%, without a facing
translation.}


\end{courseunit}

\begin{courseunit}{III.8}
\chapter{Prefaces and the Roman Canon}

\begin{learningobjectives}
  \item hold a long periodic sentence in memory by identifying its frame;
  \item read formulaic accusative--infinitive and \Latin{uti}-clauses;
  \item resolve the Canon's adjective chains and relative clauses;
  \item distinguish grammatical familiarity from theological paraphrase.
\end{learningobjectives}

The Preface begins with a stable frame and inserts a seasonal reason for praise.
The Canon is highly formulaic but syntactically dense. Because both become
familiar by hearing, use random-entry drills: begin at a clause chosen by page
and line, identify its governor, then reconnect it to the whole.

The \Term{Preface} is the thanksgiving prayer introduced by the Preface
dialogue and concluded by the angelic praise leading into the Sanctus. The
\Term{Roman Canon} is the fixed Eucharistic Prayer of the 1962 Roman Mass,
beginning \Latin{Te igitur} and ending with its doxology. A text is
\Term{formulaic} when stable wording or structures recur as recognized
formulas; formulaic does not mean syntactically unanalyzable. A
\Term{periodic sentence}, recalling the Collect definition, keeps a governing
construction incomplete while dependent members accumulate toward its
resolution. Its \Term{frame} is the stable outer construction into which those
members fit.

\section{The Preface frame}

When this chapter asks for the \Term{frame}, mark the minimum words that remain
syntactically incomplete until the thanksgiving action is supplied. Material
inserted within that frame explains why the action is fitting; it does not
become a second main sentence merely because it is long.

\begin{quote}
\Latin{Vere dignum et iustum est, æquum et salutare, nos tibi semper et
ubique gratias agere: Domine sancte, Pater omnipotens, æterne Deus.}

\emph{Source:} \Missal, \emph{Præfatio communis} and common Preface opening.
\end{quote}

The delayed infinitive clause is the true subject of impersonal \Latin{est}:
\Latin{nos} is accusative subject of \Latin{agere}, \Latin{gratias} its object,
and \Latin{tibi} dative. Four neuter predicate adjectives evaluate that action.
The three vocative titles are appositional; \Latin{æterne} is vocative masculine
despite its adverb-like appearance.

The Christmas insertion continues:
\begin{quote}
\Latin{Quia per incarnati Verbi mysterium nova mentis nostræ oculis lux tuæ
claritatis infulsit: ut, dum visibiliter Deum cognoscimus, per hunc in
invisibilium amorem rapiamur.}

\emph{Source:} \Missal, \emph{Præfatio de Nativitate}.
\end{quote}

The subject of \Latin{infulsit} is \Latin{nova ... lux}; dative
\Latin{mentis nostræ oculis} names the eyes of the mind to which it shone.
\Latin{tuæ claritatis} defines the light. The purpose/result clause contrasts
adverb \Latin{visibiliter} with substantival genitive plural
\Latin{invisibilium}. \Latin{per hunc} grammatically resumes masculine
\Latin{Deum}, the God known visibly in the mystery of the incarnate Word; it
cannot agree directly with neuter \Latin{Verbi}. Passive \Latin{rapiamur}
governs motion \Latin{in ... amorem}.

\section{The opening petition of the Canon}

\begin{quote}
\Latin{Te igitur, clementissime Pater, per Iesum Christum Filium tuum Dominum
nostrum, supplices rogamus ac petimus, uti accepta habeas et benedicas hæc
dona, hæc munera, hæc sancta sacrificia illibata.}

\emph{Source:} \Missal, Canon Missæ, \emph{Te igitur}.
\end{quote}

The fronted \Latin{te} is the object of both \Latin{rogamus} and
\Latin{petimus}; predicative \Latin{supplices} describes the implicit
\Latin{nos}. The \Latin{uti}-clause states the request. \Latin{accepta habeas}
is a resultative idiom: hold or regard the gifts as accepted.
\Latin{benedicas} shares the threefold object, whose accumulation is rhetorical
rather than three unrelated offerings. \Latin{Filium tuum Dominum nostrum} is
in apposition to \Latin{Iesum Christum} after \Latin{per}.

\section{Predicate adjectives in the \emph{Quam oblationem}}

\begin{quote}
\Latin{Quam oblationem tu, Deus, in omnibus, quæsumus, benedictam, adscriptam,
ratam, rationabilem, acceptabilemque facere digneris: ut nobis Corpus et Sanguis
fiat dilectissimi Filii tui Domini nostri Iesu Christi.}

\emph{Source:} \Missal, Canon Missæ, \emph{Quam oblationem}.
\end{quote}

Relative \Latin{quam oblationem} is the accusative object of \Latin{facere}.
Five predicate accusatives state what the offering is asked to be made.
Deponent \Latin{digneris} governs \Latin{facere}; \Latin{quæsumus} remains
parenthetical. In the result clause, singular \Latin{fiat} with \Latin{Corpus
et Sanguis} may reflect proximity agreement or a unit construal; morphology
does not decide between them. \Latin{nobis} is dative of advantage. The long
genitive chain belongs to both body and blood.

\section{Command plus passive infinitive}

\begin{quote}
\Latin{Supplices te rogamus, omnipotens Deus: iube hæc perferri per manus
sancti Angeli tui in sublime altare tuum, in conspectu divinæ maiestatis tuæ.}

\emph{Source:} \Missal, Canon Missæ, \emph{Supplices te rogamus}.
\end{quote}

Again \Latin{supplices} is predicative with implicit \Latin{nos}, and
\Latin{te} is the object of \Latin{rogamus}. Imperative \Latin{iube} governs
passive infinitive \Latin{perferri}; neuter plural \Latin{hæc} is its
accusative subject. The two \Latin{in} phrases express destination and presence
with different cases: accusative \Latin{in sublime altare}, ablative
\Latin{in conspectu}.

\section{Canon families}

Build a ledger around \Latin{memoria, memento, commemoro}; \Latin{offero,
oblatio, munus}; \Latin{accipio, acceptus, acceptabilis}; \Latin{benedico,
benedictus}; \Latin{communico, communicantes, communio}; and \Latin{supplex,
supplices}. Record the construction with each word, not only an English gloss.

\MasteryGate{you can enter the Preface or Canon at ten randomly selected
clauses, name the governing verb and reconnect each clause to its larger period;
you score at least 90\% on thirty sampled forms and antecedents and give a direct
sense of the complete Roman Canon with no facing translation. Memorized flow
without random-entry analysis does not pass.}


\end{courseunit}

\begin{courseunit}{III.9}
\chapter{Reading the Temporal Cycle}

\begin{learningobjectives}
  \item use season and genre as controlled expectations rather than guesses;
  \item recognize temporal deixis and recurrent seasonal vocabulary;
  \item read a whole formulary across chant, oration, Epistle, and Gospel registers.
\end{learningobjectives}

The \emph{Proprium de Tempore} teaches by recurrence. Advent asks for coming;
Lent for purification and perseverance; Paschaltide speaks with perfects and
alleluias of accomplished victory; Pentecost emphasizes indwelling and gift.
Such expectations suggest a semantic field, but endings still decide syntax.
Never translate \Latin{adventus} merely as ``Advent'' when the sentence requires
the concrete coming.

The \Term{Temporal cycle} is the sequence of Masses organized around Sundays,
seasons, and movable events in the Lord's mysteries. A Mass's \Term{Proper}
comprises texts assigned to that occasion rather than fixed across Masses; a
\Term{formulary}, as defined at the opening of this volume, gathers those
assigned texts into one Mass. A \Term{register} is a recognizable variety of
language associated with a function or textual setting---for example, prayer,
narrative, rubric, or psalmic poetry---not a ranking of good and bad Latin.

\Term{Deixis} is linguistic pointing from the text's situation: words such as
\Latin{hodie}, \Latin{hic}, and \Latin{ventura} locate time, place, or
participants relative to the act of speaking. A \Term{semantic field} is a
group of words associated by an area of meaning; it can guide expectation but
does not determine any word's morphology or contextual sense. A \Term{lexeme}
is a vocabulary item considered across all its inflected forms, whereas a
lemma is the dictionary form used to list it.

\section{Advent: descent and germination}

\begin{quote}
\Latin{Rorate cæli desuper, et nubes pluant iustum: aperiatur terra, et
germinet Salvatorem.}

\emph{Source:} \Missal, Dominica IV Adventus, Introitus (Is.~45:8).
\end{quote}

\Latin{rorate} is an imperative addressed to plural \Latin{cæli}; adverb
\Latin{desuper} gives direction. Jussive \Latin{pluant, aperiatur}, and
\Latin{germinet} form a sequence. Substantival accusative \Latin{iustum} is what
the clouds rain; \Latin{Salvatorem} is the object brought forth by the earth.
The capitalization is devotional typography, not a different case.

\section{Lent: practice directed toward action}

\begin{quote}
\Latin{Deus, qui Ecclesiam tuam annua quadragesimali observatione purificas:
præsta familiæ tuæ; ut, quod a te obtinere abstinendo nititur, hoc bonis
operibus exsequatur.}

\emph{Source:} \Missal, Dominica I in Quadragesima, Collecta.
\end{quote}

The relative clause supplies the seasonal reason for petition. Ablative
\Latin{annua quadragesimali observatione} is means. Dative
\Latin{familiæ tuæ} receives the grant whose content follows. Correlative
\Latin{quod ... hoc} marks the same object: what the household strives to
obtain by fasting, let it carry out in good works. \Latin{abstinendo} is an
ablative gerund; deponent \Latin{nititur} takes complementary \Latin{obtinere};
\Latin{exsequatur} is deponent subjunctive.

\section{Easter: the startling perfect}

\begin{quote}
\Latin{Resurrexi, et adhuc tecum sum, alleluia: posuisti super me manum tuam,
alleluia: mirabilis facta est scientia tua, alleluia, alleluia.}

\emph{Source:} \Missal, Dominica Resurrectionis, Introitus (Ps.~138:18, 5--6).
\end{quote}

In its Easter placement, traditional liturgical reception hears the risen
Christ as speaker. Perfect \Latin{resurrexi} stands beside present \Latin{sum};
\Latin{tecum} combines \Latin{cum} with the pronoun. In \Latin{mirabilis facta
est scientia tua}, \Latin{mirabilis} is predicate nominative and \Latin{scientia
tua} subject. That speaker identification belongs to placement and reception;
morphology alone does not prove it.

\section{Pentecost: a relative containing all things}

\begin{quote}
\Latin{Spiritus Domini replevit orbem terrarum, alleluia: et hoc quod continet
omnia, scientiam habet vocis, alleluia.}

\emph{Source:} \Missal, Dominica Pentecostes, Introitus (Sap.~1:7).
\end{quote}

\Latin{Spiritus} is nominative singular and \Latin{Domini} genitive;
\Latin{orbem terrarum} is object plus defining genitive. Demonstrative
\Latin{hoc} is the subject of \Latin{habet} and is expanded by
\Latin{quod continet omnia}; \Latin{scientiam} is object and \Latin{vocis}
genitive of what is known.

\section{The seasonal ledger}

For each formulary record four columns: repeated lexemes, time words, governing
petitions, and biblical allusions. Especially notice \Latin{hodie, annuus,
ventura, præsens, æternus}; \Latin{adventus, exspecto, venio};
\Latin{ieiunium, abstinentia, purifico}; \Latin{resurgo, renovatio, vita}; and
\Latin{Spiritus, donum, infusio, repleo}. A predicted family earns no parse
credit until the actual form is identified.

\MasteryGate{using a randomly assigned Sunday from each of Advent, Lent,
Paschaltide, and the season after Pentecost, you read every Proper text without
a facing translation, achieve at least 90\% clause and form accuracy on a
sampled audit, and produce a one-page genre map showing how the formulary's
chants, prayers, and readings cohere.}


\end{courseunit}

\begin{courseunit}{III.10}
\chapter{The Sanctoral Cycle and the Commons}

\begin{learningobjectives}
  \item distinguish a saint's Proper from a supplied Common;
  \item adapt placeholder formulas without losing agreement;
  \item recognize recurrent martyr, confessor, virgin, and Marian vocabulary;
  \item avoid mistaking remembered formulas for the text actually assigned.
\end{learningobjectives}

The \emph{Proprium Sanctorum} may provide a complete formulary or direct the
reader to a \emph{Commune}. A Common is not a single generic Mass: it is a set
of formularies and options classified by saint and office. Before reading, write
the assignment exactly: one martyr or several, pontiff or non-pontiff,
confessor, doctor, virgin, holy woman, dedication, or the Blessed Virgin Mary.
Then verify which elements remain proper.

The \Term{Sanctoral cycle} is the calendar sequence of saints' feasts and other
fixed-date observances. A saint's \Term{Proper} supplies text particular to
that observance; a \Term{Common} supplies one or more reusable formularies for
a classified kind of saint. A \Term{placeholder} such as \Latin{N.} marks a
slot to be filled from the assignment. Filling it is a grammatical operation:
the supplied name must take the case required by the surrounding formula.

\section{The grammatical placeholder}

\begin{quote}
\Latin{Deus, qui nos beati N. Martyris tui atque Pontificis annua solemnitate lætificas:
concede propitius; ut, cuius natalitia colimus, de eiusdem etiam protectione
gaudeamus.}

\emph{Source:} \Missal, Commune unius Martyris extra tempus paschale, Pro
Martyre Pontifice, Altera Missa, Oratio.
\end{quote}

\Latin{beati N. Martyris tui atque Pontificis} is a genitive phrase depending on
\Latin{solemnitate}. The abbreviation \Latin{N.} must be expanded with the
saint's name in the case demanded by the formula, not copied mechanically in
nominative. Causative \Latin{lætificas} takes \Latin{nos}. Correlative
\Latin{cuius ... eiusdem} refers to the martyr: the feast day is celebrated and
protection from the same person is enjoyed. \Latin{de ... protectione} belongs
with \Latin{gaudeamus}.

\section{A Common Gospel command}

\begin{quote}
\Latin{Si quis vult post me venire, abneget semetipsum, et tollat crucem suam,
et sequatur me.}

\emph{Source:} \Missal, Commune unius Martyris extra tempus paschale, Pro
Martyre Pontifice, Altera Missa, Evangelium (Mt.~16:24--27).
\end{quote}

Indefinite \Latin{si quis} introduces a present open condition. Complementary
\Latin{venire} depends on \Latin{vult}; \Latin{post me} expresses following.
The apodosis has three coordinated jussive subjunctives. Intensive reflexive
\Latin{semetipsum} is the object of \Latin{abneget}; \Latin{suam} refers to the
same individual subject, while \Latin{me} refers to the speaker.

\section{Wisdom on the lips}

\begin{quote}
\Latin{Os iusti meditabitur sapientiam, et lingua eius loquetur iudicium: lex
Dei eius in corde ipsius.}

\emph{Source:} \Missal, Commune Confessoris non Pontificis, Introitus
(Ps.~36:30--31).
\end{quote}

\Latin{os} and \Latin{lingua} are subjects of future deponent
\Latin{meditabitur} and future \Latin{loquetur}. The Psalm's verbs take
accusatives \Latin{sapientiam} and \Latin{iudicium}. The final clause is
verbless; a bracketed \Latin{[est]} is an interpretive dependency or English
paraphrase, not an expressed or grammatically required copula. Distinguish possessive
\Latin{eius} from intensive/genitive \Latin{ipsius}; both refer to the just
person, while \Latin{Dei} identifies the law.

\section{The Marian Common}

\begin{quote}
\Latin{Salve, sancta Parens, enixa puerpera Regem: qui cælum terramque regit
in sæcula sæculorum.}

\emph{Source:} \Missal, Commune Festorum Beatæ Mariæ Virginis, Introitus
(Sedulius).
\end{quote}

\Latin{Parens} and \Latin{puerpera} are vocative in function; participle
\Latin{enixa} agrees with the addressed mother and governs accusative
\Latin{Regem}. Relative \Latin{qui} refers to the King, not to the mother, as
masculine gender and \Latin{regit} confirm. Enclitic \Latin{-que} joins
\Latin{cælum} and \Latin{terram}; \Latin{in sæcula sæculorum} is biblical
genitive intensification, a formula of unbounded duration.

\section{Name and agreement ledger}

For each saint write the nominative and the genitive actually used by the
Missal. Then write titles in the required gender and number: \Latin{Martyris,
Confessoris, Pontificis, Virginis, Viduæ}. When a Common changes from one saint
to several, check every pronoun, participle, relative, and verb; do not merely
pluralize the name.

\MasteryGate{given four unseen sanctoral assignments, you correctly locate
Proper and Common components, expand every N. in the demanded case, adapt a
singular formula to a supplied plural subject without agreement errors, and
read each complete formulary with at least 90\% sampled parsing accuracy and no
facing translation.}


\end{courseunit}

\begin{courseunit}{III.11}
\chapter{Requiem and Holy Week}

\begin{learningobjectives}
  \item track singular, plural, and gender changes in Masses for the dead;
  \item read the compressed eschatological diction of the Requiem chants;
  \item recognize Holy Week's changes of register and speaker;
  \item give safety priority to the actual 1962 text rather than a remembered text from another edition.
\end{learningobjectives}

These corpora punish assumption. Requiem formularies vary with one deceased or
many and with grammatical gender. Holy Week contains blessings, narrative,
solemn prayer, acclamation, reproach, and the Paschal proclamation. Always read
the heading and the printed 1962 formula; texts remembered from earlier or later
editions may differ.

A \Term{Requiem Mass} is a Mass formulary for the dead; the label covers
several 1962 formularies whose wording changes with occasion and intention.
An \Term{acclamation} is a short public utterance of praise, assent, or
recognition; a \Term{reproach} is direct address framed as an accusation or
grieved question. A \Term{register change} occurs when the text's discourse
job and characteristic language change---for example, from narrative to
command or lament---even inside the same day's rites. Recall that
\Term{deixis} locates persons, time, or place relative to the utterance; in a
Requiem, resolving \Latin{ei/eis}, \Latin{hodie}, and similar forms therefore
requires the exact formulary and intention.

\section{Requiem deixis and agreement}

\begin{quote}
\Latin{Requiem æternam dona eis, Domine: et lux perpetua luceat eis. Te decet
hymnus, Deus, in Sion, et tibi reddetur votum in Ierusalem.}

\emph{Source:} \Missal, Missæ Defunctorum, Introitus
(4~Esd.~2:34--35; Ps.~64:2).
\end{quote}

Imperative \Latin{dona} takes object \Latin{requiem æternam} and dative
\Latin{eis}; jussive \Latin{luceat} takes subject \Latin{lux perpetua} and the
same dative. In \Latin{te decet hymnus}, accusative \Latin{te} is the person
whom the hymn befits and \Latin{hymnus} is subject. Passive future
\Latin{reddetur} takes \Latin{votum} as subject and \Latin{tibi} as recipient.

For one deceased a collect may say:
\begin{quote}
\Latin{Deus, cui proprium est misereri semper et parcere: te supplices exoramus
pro anima famuli tui N., quam hodie de hoc sæculo migrare iussisti.}

\emph{Source:} \Missal, Missæ Defunctorum, in die obitus seu depositionis,
Collecta.
\end{quote}
\Latin{cui} is dative of possession with \Latin{proprium est}; the subject is
the paired infinitive phrase \Latin{misereri ... et parcere}. Fronted
\Latin{te} is object of \Latin{exoramus}; \Latin{supplices} describes the
implicit petitioners. Relative \Latin{quam} refers to \Latin{animam} and is the
accusative subject of \Latin{migrare} after \Latin{iussisti}. Substituting a
woman changes \Latin{famuli} to \Latin{famulæ}; a plural formulary changes much
more than that.

\section{The sequence as compressed judgment}

\begin{quote}
\Latin{Dies iræ, dies illa, solvet sæclum in favilla: teste David cum
Sibylla.}

\emph{Source:} \Missal, Missæ Defunctorum, Sequentia \emph{Dies iræ}.
\end{quote}

Although \Latin{dies} is often masculine, it can be feminine when it denotes
an appointed or fixed day; \Latin{illa} therefore agrees regularly with
\Latin{dies} here. Future \Latin{solvet} takes \Latin{sæclum} as object;
\Latin{in favilla} expresses result. Ablative absolute
\Latin{teste David} means David bearing witness; \Latin{cum Sibylla} adds the
second witness. Poetry may preserve constructions that prose drills would not
choose.

\section{Holy Thursday and obligation}

\begin{quote}
\Latin{Nos autem gloriari oportet in cruce Domini nostri Iesu Christi: in quo
est salus, vita et resurrectio nostra: per quem salvati et liberati sumus.}

\emph{Source:} \Missal, Feria V in Cena Domini, Introitus (Gal.~6:14).
\end{quote}

Impersonal \Latin{oportet} governs infinitive \Latin{gloriari}; \Latin{nos} is
the infinitive's accusative subject. Relative \Latin{quo} refers to
\Latin{Domino nostro Iesu Christo} as the person in whom salvation is found,
while \Latin{quem} after \Latin{per} names the agent through whom the plural
subject has been saved and freed. The preverbal singular \Latin{est} with
\Latin{salus, vita et resurrectio nostra} may show proximity agreement with
the first noun or treat the series as a unit; morphology alone does not decide.

\section{Reproach and acclamation}

\begin{quote}
\Latin{Popule meus, quid feci tibi? aut in quo contristavi te? responde mihi.}

\emph{Source:} \Missal, Feria VI in Parasceve, Improperia.
\end{quote}

\Latin{Popule meus} is vocative. Interrogative \Latin{quid} is object of
\Latin{feci}, dative \Latin{tibi} marks the affected person, and
\Latin{in quo} asks in what respect. In \Latin{contristavi te}, \Latin{te} is
the direct object. The command \Latin{responde} takes dative \Latin{mihi}.

The veneration acclamation says:
\begin{quote}
\Latin{Ecce lignum Crucis, in quo salus mundi pependit. Venite, adoremus.}

\emph{Source:} \Missal, Feria VI in Parasceve, Adoratio Crucis.
\end{quote}
\Latin{ecce} presents \Latin{lignum}; because this neuter form is identical in
the nominative and accusative, its case cannot be proved from the ending alone
(\Latin{ecce} admits both constructions). \Latin{in quo} is an ablative
relative, \Latin{salus} the subject, and \Latin{mundi} a genitive. Plural
imperative \Latin{venite} is followed by hortatory \Latin{adoremus}.

\section{The Paschal proclamation}

\begin{quote}
\Latin{Exsultet iam angelica turba cælorum: exsultent divina mysteria: et pro
tanti Regis victoria tuba insonet salutaris.}

\emph{Source:} \Missal, Sabbato Sancto, Præconium paschale.
\end{quote}

Three jussive subjunctives have three subjects: \Latin{turba},
\Latin{mysteria}, and \Latin{tuba}. The apparent singular/plural surprise in
\Latin{angelica turba cælorum} resolves when \Latin{turba} is recognized as a
collective singular. \Latin{pro ... victoria} gives the occasion, with
\Latin{tanti Regis} dependent on \Latin{victoria}; \Latin{salutaris} agrees
with \Latin{tuba}.

\MasteryGate{you can adapt a Requiem prayer accurately for one man, one woman,
and several deceased persons; parse an unseen stanza of \emph{Dies iræ}; and
read one assigned Holy Week office section while correctly marking every change
of speaker and register, all without a facing translation and at 90\% sampled
form accuracy.}


\end{courseunit}

\begin{courseunit}{III.12}
\chapter{Full-Mass Unseen Capstones}

\begin{learningobjectives}
  \item execute a complete formulary reading without leaning on remembered English;
  \item keep a disciplined lookup and error ledger;
  \item demonstrate direct comprehension through Latin answers and grammatical evidence;
  \item decide honestly whether to advance, remediate a genre, or repeat a gate.
\end{learningobjectives}

The capstone is not a race and not a performance of vague spiritual familiarity.
It is an audit trail. Use a clean copy of the \Missal, blank paper, a Latin
dictionary, and no parallel translation. A grammar may be consulted only after
the timed first pass and must be logged like a dictionary lookup.

A \Term{capstone} is an integrative assessment requiring earlier skills to
operate together on a substantial new task. An \Term{audit trail} is the
written record that makes the process checkable: marks, proposed parses,
lookups, revisions, and classified errors. \Term{Direct comprehension} means
that the learner can follow agents, actions, relations, and discourse movement
from the Latin itself. A \Term{literal sense} is a restrained account that
preserves those grammatical relations; it need not mimic Latin word order or
replace each Latin word mechanically with one English word.

\section{The complete-Mass protocol}

\begin{enumerate}
  \item \textbf{Inventory, 5 minutes.} List the formulary and every genre it
    contains. Identify directions to a Common or another Mass.
  \item \textbf{Cold read.} Read all Proper texts once without a lookup. Mark
    finite verbs, clause boundaries, speakers, and only the most likely heads.
  \item \textbf{Analytical read.} Complete dependencies, prose-order lines for
    compressed passages, and the ellipsis ledger. Add the fixed Ordo and Canon
    by random-entry prompts rather than reciting from the beginning.
  \item \textbf{Controlled lookup.} Look up only marked content words after
    proposing lemma, morphology, and function. Record every lookup.
  \item \textbf{Direct-comprehension proof.} Answer twelve questions in simple
    Latin: agents, actions, recipients, causes, purposes, contrasts, and temporal
    sequence. Then give a restrained literal written sense.
  \item \textbf{Audit.} Check against the analytical key or a competent reader.
    Classify every error as form, construction, vocabulary, reference, ellipsis,
    or unsupported inference.
\end{enumerate}

For that classification, a \Term{form error} misidentifies or produces
morphology; a \Term{construction error} misstates a syntactic dependency; a
\Term{vocabulary error} selects the wrong contextual sense; a \Term{reference
error} assigns a pronoun or modifier to the wrong referent; an \Term{ellipsis
error} supplies or omits wording without grammatical warrant; and an
\Term{unsupported inference} asserts more than the textual evidence permits.

\section{Dictionary thresholds}

The target is not zero dictionary use. Mature readers use dictionaries to
distinguish senses. The target is no dictionary dependence for grammatical
structure. On a first capstone, permit ten new content-word lookups per five
hundred words. On the second, eight; on the final two, five. A repeated lookup
of a word already in the ledger counts twice. Proper names do not count, but
failing to identify their case does.

\section{Assigned capstones}

Use texts not excerpted in this volume:
\begin{enumerate}
  \item \textbf{Capstone A---Temporal prose:} Dominica XII post Pentecosten.
    Emphasize the Collect, 2~Cor.~3:4--9, and Lc.~10:23--37.
  \item \textbf{Capstone B---Temporal contrast:} Dominica XXIII post
    Pentecosten. Read all chants and prayers as well as Phil.~3:17--21; 4:1--3
    and Mt.~9:18--26.
  \item \textbf{Capstone C---Common and supplied name:} a locally available
    feast using the Commune unius Martyris non Pontificis. Record every supplied
    element and agreement change.
  \item \textbf{Capstone D---Requiem:} Missa quotidiana Defunctorum, adapting
    the prayers to a plural intention supplied by the assessor.
\end{enumerate}

These assignments test breadth but not every page of the Missal. Long-term
independence grows by sampling the complete annual cycle: after each review
cycle, choose another previously unread Sunday or feast formulary and follow it
with an error-ledger review. Pace that work by mastery rather than by weeks.
Familiar Ordinaries and Commons must not be allowed to inflate the unseen
score.

\section{Measurable final gate}

Score each capstone in five equal domains:
\begin{center}
\begin{tabularx}{\linewidth}{@{}P{1.45in}Y@{}}
\toprule
\textbf{Domain} & \textbf{Passing evidence} \\
\midrule
Morphology & At least 90\% correct on forty randomly sampled forms. \\
Clause structure & At least 90\% of finite clauses correctly bounded and attached. \\
Reference/ellipsis & At least 90\% of sampled pronouns, participles, and supplies justified. \\
Direct comprehension & At least 10 of 12 Latin questions substantively correct; literal sense preserves agent, action, negation, and purpose. \\
Dictionary discipline & Lookup ceiling observed; every lookup begins with a proposed lemma and function. \\
\bottomrule
\end{tabularx}
\end{center}

Passing requires every domain at threshold on two consecutive unseen capstones.
An average cannot conceal a failed syntax or comprehension domain. If a domain
fails, return to the matching genre chapter, complete two unused worksheet
variants, and attempt a different unseen formulary.

\MasteryGate{you pass two consecutive full-Mass unseen capstones in every
domain, then read a third complete formulary at normal silent reading pace with
no facing translation and no more than five content-word lookups per five
hundred words. You can also write a 120-word controlled Latin summary whose
agents, tenses, agreements, and subordinate relations contain no more than five
substantive errors after one self-correction pass.}
\end{courseunit}
